\section*{Chapter \ref{ch:b2:current-potential-curves} --- Current--Potential Curves}

\begin{solution}{exo:b2:current-potential-curves:1}
Zinc is oxidised: \omterm{def:b2:current-potential-curves:ie-curve}{anodic current}, positive. Copper ions are reduced at the
copper electrode: \omterm{def:b2:current-potential-curves:ie-curve}{cathodic current}, negative. The two currents are equal in
magnitude, as the same charge flows through the circuit.
\end{solution}

\begin{solution}{exo:b2:current-potential-curves:2}
$|i_{\lim}| = 96\,485 \times 0.50 \times 10^{-4} \times 1.6 \times 10^{-9} \times
2.0/(25 \times 10^{-6}) = \qty{6.2e-4}{A}$, that is \qty{0.62}{mA}.
\end{solution}

\begin{solution}{exo:b2:current-potential-curves:3}
A \omterm{def:b2:current-potential-curves:fast-slow}{fast system} crosses the potential axis steeply at its Nernst potential; a \omterm{def:b2:current-potential-curves:fast-slow}{slow
system} shows a flat stretch of zero current around it, its branches starting
only after an \omterm{def:b2:current-potential-curves:overpotential}{overpotential}. On platinum: \ce{Fe^3+}/\ce{Fe^2+} (fast);
\ce{O2}/\ce{H2O} (slow).
\end{solution}

\begin{solution}{exo:b2:current-potential-curves:4}
At the half-wave potential, equal to $E^\circ = \qty{0.77}{V}$ when the two forms
diffuse alike.
\end{solution}

\begin{solution}{exo:b2:current-potential-curves:5}
Near \qty{0.77}{V} a cathodic wave of \ce{Fe^3+} (fast), reaching a plateau
proportional to its concentration. Near \qty{0.34}{V} a second wave, copper being
deposited; its plateau adds to the first, about twice as high for an equal
concentration and similar diffusion coefficients, since $n = 2$. Then, near
\qty{0}{V} at pH 0 on platinum (lower once the electrode is coated with copper),
the steep wall of the reduction of \ce{H+}.
\end{solution}

\begin{solution}{exo:b2:current-potential-curves:6}
At \qty{0.50}{V}: only \ce{Fe^3+} is reduced to \ce{Fe^2+}, at its \omterm{def:b2:current-potential-curves:limiting-current}{limiting current}.
At \qty{0.10}{V}: \ce{Fe^3+} is reduced and copper is deposited at the same time,
each at its \omterm{def:b2:current-potential-curves:limiting-current}{limiting current}; the total is the sum of the two plateaus.
\end{solution}

\begin{solution}{exo:b2:current-potential-curves:7}
pH 0: \qty{0}{V} and \qty{1.23}{V}; pH 14: \qty{-0.83}{V} and \qty{0.40}{V}. On
platinum the window keeps its width and moves down by \qty{59}{mV} per pH unit,
each wall following its couple, the anodic one still shifted up by the oxygen
\omterm{def:b2:current-potential-curves:overpotential}{overpotential}.
\end{solution}

\begin{solution}{exo:b2:current-potential-curves:8}
At \qty{1.0}{V}, \ce{Fe^2+} is oxidised on its plateau (\omterm{def:b2:current-potential-curves:ie-curve}{anodic current}
proportional to $[\ce{Fe^2+}]$) and any \ce{Ce^4+} is reduced on its plateau
(\omterm{def:b2:current-potential-curves:ie-curve}{cathodic current} proportional to $[\ce{Ce^4+}]$). Before the equivalence the
\omterm{def:b2:current-potential-curves:ie-curve}{anodic current} falls linearly with the volume added; after it the \omterm{def:b2:current-potential-curves:ie-curve}{cathodic
current} grows linearly. The two straight lines cross zero at the equivalence.
\end{solution}

\begin{solution}{exo:b2:current-potential-curves:9}
The plateaus double ($|i_{\lim}| \propto 1/\delta$). The half-wave potential does
not change, the layer being the same for both forms; nor does the zero-current
potential of a \omterm{def:b2:current-potential-curves:fast-slow}{fast system}, which is the Nernst potential.
\end{solution}

\begin{solution}{exo:b2:current-potential-curves:10}
As in the chapter with $i_{\lim,c} = 0$: $E = E_{1/2} + (RT/nF)\ln[i/(i_{\lim} -
i)]$. In decimal logarithms, $E = E_{1/2} + (RT\ln 10/nF)\log[i/(i_{\lim} - i)]$,
a straight line of slope $0.0592/n$ volt at \qty{298}{K}, which crosses the
$E$ axis at $E_{1/2}$.
\end{solution}

\begin{solution}{exo:b2:current-potential-curves:11}
Alone, the zinc settles at the potential where its \omterm{def:b2:current-potential-curves:ie-curve}{anodic current} (fast
oxidation of zinc) equals the \omterm{def:b2:current-potential-curves:ie-curve}{cathodic current} of \ce{H+} reduction on zinc,
which is small because that reduction is slow on zinc: the dissolution is
slow. Touching platinum, the electrons released by the zinc can reduce \ce{H+}
on the platinum, where the reduction is fast: the \omterm{def:b2:current-potential-curves:ie-curve}{cathodic current} at a given
potential is much larger, the balance is reached at a higher current, and the
zinc dissolves faster while hydrogen forms on the platinum.
\end{solution}

\begin{solution}{exo:b2:current-potential-curves:12}
On platinum the oxidation of water starts near \qty{1.6}{V} at pH 0: before
\qty{2.0}{V} the current would be spent making dioxygen. An electrode on which
the oxidation of water needs a larger \omterm{def:b2:current-potential-curves:overpotential}{overpotential} pushes the wall beyond
\qty{2.0}{V}, so that the species can be oxidised inside the window.
\end{solution}

\begin{solution}{pb:b2:current-potential-curves:1}
\textbf{1.} \ce{O2 + 2H2O + 4e- -> 4OH-}.
\textbf{2.} \ce{Ag + Cl- -> AgCl + e-} (four times); overall \ce{O2 + 2H2O + 4Ag +
4Cl- -> 4AgCl + 4OH-}.
\textbf{3.} $E = 1.229 - 0.0592 \times 7 + (0.0592/4)\log 0.21 = \qty{0.80}{V}$.
\textbf{4.} $E^\circ(\ce{AgCl}/\ce{Ag}) = \qty{0.22}{V}$ with chloride of activity 1:
$-0.60 + 0.22 = \qty{-0.38}{V}$ on the hydrogen scale.
\textbf{5.} The reduction of \ce{O2} is a \omterm{def:b2:current-potential-curves:fast-slow}{slow system}: below its Nernst potential
an \omterm{def:b2:current-potential-curves:overpotential}{overpotential} of several tenths of a volt is needed before the current
grows, and more to reach the plateau.
\textbf{6.} On the plateau the current depends only on the supply of \ce{O2},
not on small drifts of the potential: it measures the concentration.
\textbf{7.} The reduction of water to dihydrogen (the cathodic wall), which would
add a current unrelated to oxygen.
\textbf{8.} Linear across the membrane, from $c$ on the sample side to zero at
the cathode.
\textbf{9.} $j = D_m c/L$ per unit area.
\textbf{10.} $i = -4FAD_m c/L$ (cathodic), $|i| = 4FAD_mc/L$.
\textbf{11.} The membrane is the slowest step: diffusion across it is much
slower than in the water outside, so the membrane plays the part of the
\omterm{def:b2:current-potential-curves:limiting-current}{diffusion layer}.
\textbf{12.} The concentration gradient lies inside the membrane, whose thickness
does not depend on the stirring (provided the water outside is renewed);
a deposit on the membrane adds thickness and lowers the current.
\textbf{13.} $0.2095 \times (101.325 - 3.17) = \qty{20.56}{kPa}$.
\textbf{14.} $c = 1.3 \times 10^{-5} \times 20\,560 = \qty{0.27}{mol/m^3}$ ($0.267$).
\textbf{15.} $0.267 \times 32.00 = \qty{8.6}{mg/L}$ (\unit{g/m^3}).
\textbf{16.} $|i| = 4 \times 96\,485 \times 1.0 \times 10^{-5} \times 1.0 \times 10^{-10}
\times 0.267/(25 \times 10^{-6}) = \qty{4.1}{\micro A}$.
\textbf{17.} $4.05/8.55 = \qty{0.474}{\micro A}$ per \unit{mg/L}.
\textbf{18.} The membrane thickness and its diffusion coefficient are not known
precisely and change with age and temperature; a measurement in a known
solution includes them all.
\textbf{19.} Zero (in practice a small residual current, subtracted).
\textbf{20.} $2.80/0.474 = \qty{5.9}{mg/L}$.
\textbf{21.} $2.80/4.05 = 69~\%$ of saturation.
\textbf{22.} Oxygen consumed by the decay of organic matter (sewage, dead plants)
faster than it is supplied by the air and by photosynthesis.
\textbf{23.} $2.80 \times 10^{-6}/(4 \times 96\,485) \times 3600 = \qty{2.6e-8}{mol}$
per hour: negligible against the oxygen of even a litre of water
(\qty{1.8e-4}{mol}).
\textbf{24.} Both the solubility of \ce{O2} (Henry's constant) and the diffusion
across the membrane change with temperature: the sensitivity found at
\qty{25}{\degreeCelsius} no longer applies.
\textbf{25.} The sample holds $\boldsymbol{\approx \qty{5.9}{mg/L}}$ of dissolved
oxygen.
\end{solution}
